\documentclass[11pt,a4paper]{article}

\usepackage[utf8]{inputenc}
\usepackage[T1]{fontenc}
\usepackage[french]{babel}
\usepackage[a4paper,margin=2.2cm]{geometry}
\usepackage{amsmath,amssymb,amsthm}
\usepackage{lmodern}
\usepackage{enumitem}
\usepackage{xcolor}
\usepackage{fancyhdr}
\usepackage{lastpage}
\usepackage{titlesec}
\usepackage{tcolorbox}
\usepackage{array}
\usepackage{tabularx}
\usepackage{multicol}
\usepackage{hyperref}

\hypersetup{
    colorlinks=true,
    linkcolor=black,
    urlcolor=blue
}

\pagestyle{fancy}
\fancyhf{}
\lhead{\textbf{DM de mathématiques -- Niveau Première}}
\rhead{\textbf{Nom : \hspace{4cm}}}
\cfoot{\thepage/\pageref{LastPage}}

\setlength{\parindent}{0pt}
\setlength{\parskip}{0.4em}

\definecolor{bleufonce}{RGB}{25,60,120}
\definecolor{bleuclair}{RGB}{235,242,255}
\definecolor{grisclair}{RGB}{245,245,245}

\titleformat{\section}
{\normalfont\Large\bfseries\color{bleufonce}}{\thesection.}{0.6em}{}

\titleformat{\subsection}
{\normalfont\large\bfseries\color{bleufonce}}{\thesubsection.}{0.6em}{}

\tcbset{
    colback=bleuclair,
    colframe=bleufonce,
    arc=2mm,
    boxrule=0.8pt
}

\begin{document}

\begin{center}
    {\LARGE \textbf{DM0 de mathématiques indicatif}}\\[0.4em]
    {\large Niveau Première spé}\\[0.8em]
    \rule{0.9\textwidth}{0.8pt}\\[0.6em]
\end{center}

\begin{tabularx}{\textwidth}{|>{\centering\arraybackslash}p{0.22\textwidth}|>{\centering\arraybackslash}p{0.22\textwidth}|>{\centering\arraybackslash}X|}
\hline
\textbf{Durée} & \textbf{Barème} & \textbf{Matériel autorisé} \\
\hline
1 h 30 à 2 h & 20 points & stylo, feuille \\
\hline
\end{tabularx}

\vspace{1em}

\begin{tcolorbox}[title=Consignes générales]
\begin{itemize}[leftmargin=1.2cm]
    \item Pensez bien à noter le temps que vous avez pris à faire le sujet
    \item Tout résultat \textbf{non trivial} doit être \textbf{justifié}, ou au moins pouvoir être expliqué clairement à l'oral.
    \item On attend un \textbf{raisonnement mathématique clair}, structuré et rigoureux.
    \item Merci de n'utiliser aucune aide extérieure.
    \item Lorsqu'une aide ou un rappel est donné, il peut être utilisé librement.
\end{itemize}
\end{tcolorbox}

\vspace{1em}

\begin{center}
    \renewcommand{\arraystretch}{1.4}
    \begin{tabular}{|c|c|c|c|c|c|}
        \hline
        \textbf{Exercice 1} & \textbf{Exercice 2} & \textbf{Exercice 3} & \textbf{Exercice 4} & \textbf{Exercice 5} & \textbf{Total} \\
        \hline
        12 pts & 4 pts & 4 pts & 3 pts & 4 pts & 27 pts \\
        \hline
    \end{tabular}
\end{center}

\vspace{1.2em}

% ========================= EXERCICE 1 =========================
\section*{Exercice 1 -- Fonctions exponentielles \hfill \textnormal{12 points}}

On considère la fonction $f$ définie sur $\mathbb{R}$ par :
\[
f(x)=3e^{0,5x}-2.
\]

\begin{enumerate}
    \item Calculer $f(0)$ et $f(2)$. \hfill \textbf{2 pts}

    \item Résoudre l’équation :
    \[
    f(x)=1.
    \]
    On donnera la valeur exacte de $x$, puis une valeur approchée au centième. \hfill \textbf{3 pts}

    \item Résoudre l’inéquation :
    \[
    f(x)>4.
    \]
    \hfill \textbf{3 pts}

    \item On admet que la fonction $x \mapsto e^{0,5x}$ est strictement croissante sur $\mathbb{R}$.

    En déduire le sens de variation de la fonction $f$ sur $\mathbb{R}$. \hfill \textbf{2 pts}

    \item Dans un contexte concret, on suppose que $f(x)$ représente une grandeur physique exprimée en unités, après $x$ heures.

    Déterminer au bout de combien d’heures cette grandeur dépasse $10$. \hfill \textbf{2 pts}
\end{enumerate}

% ========================= EXERCICE 2 =========================
\section*{Exercice 2 -- Démonstration avec des triangles \hfill /4}

On considère un triangle \(ABC\) isocèle en \(A\), c'est-à-dire tel que :
\[
AB=AC.
\]
On note \(M\) le milieu du segment \([BC]\).

On veut démontrer que la droite \((AM)\) est perpendiculaire à la droite \((BC)\).

\begin{tcolorbox}[title=Rappels utiles]
\begin{itemize}[leftmargin=1.2cm]
    \item Dire que \(M\) est le milieu de \([BC]\) signifie que \(M \in [BC]\) et \(BM=MC\).
    \item Si deux triangles ont leurs trois côtés respectivement égaux, alors ils sont isométriques.
    \item Des triangles isométriques ont leurs angles correspondants égaux.
    \item Deux angles adjacents dont les côtés extérieurs sont dans le prolongement l'un de l'autre sont supplémentaires.
    \item Deux angles égaux et supplémentaires mesurent chacun \(90^\circ\).
\end{itemize}
\end{tcolorbox}

\begin{enumerate}[label=\textbf{\arabic*.}, leftmargin=1.2cm]
    \item Montrer que les triangles \(ABM\) et \(ACM\) sont isométriques.
    \item En déduire que
    \[
    \widehat{BMA}=\widehat{AMC}.
    \]
    \item Expliquer pourquoi les angles \(\widehat{BMA}\) et \(\widehat{AMC}\) sont supplémentaires.
    \item Conclure que
    \[
    (AM)\perp(BC).
    \]
\end{enumerate}

\vspace{0.6em}
\textbf{Barème indicatif :} 1,5 pt ; 1 pt ; 0,5 pt ; 1 pt.

% ========================= EXERCICE 3 =========================
\section*{Exercice 3 -- Démonstration à l'aide d'un schéma \hfill /4}

\textbf{Consigne :} \emph{À l'aide d'un schéma, démontrer que la somme des angles d'un triangle est \(180^\circ\).}

\medskip

On considère un triangle quelconque \(ABC\).

\begin{enumerate}[label=\textbf{\arabic*.}, leftmargin=1.2cm]
    \item Faire un schéma clair du triangle \(ABC\).
    \item Par le point \(A\), tracer mentalement ou sur le schéma une droite parallèle à la droite \((BC)\).
    \item En utilisant les angles alternes-internes ou correspondants, démontrer que :
    \[
    \widehat{ABC}+\widehat{BAC}+\widehat{ACB}=180^\circ.
    \]
\end{enumerate}

\begin{tcolorbox}[title=Attendus]
La démonstration doit s'appuyer sur un schéma clair, des notations d'angles précises et un raisonnement rédigé. Il ne suffit pas d'affirmer le résultat : il faut expliquer pourquoi les angles utilisés sont égaux et pourquoi l'angle plat vaut \(180^\circ\).
\end{tcolorbox}

\vspace{0.6em}
\textbf{Barème indicatif :} schéma et notations 1 pt ; utilisation correcte du parallélisme 1,5 pt ; conclusion rédigée 1,5 pt.

% ========================= EXERCICE 4 =========================
\section*{Exercice 4 -- Logique et démonstration \hfill /3}

\begin{enumerate}[label=\textbf{\arabic*.}, leftmargin=1.2cm]
    \item Montrer que la somme de deux nombres pairs est paire.
    \item Montrer que le carré d'un nombre impair est impair.
    \item Démontrer que
    \[
    (x-4)(x+2)=0 \Longrightarrow x=4 \text{ ou } x=-2.
    \]
\end{enumerate}

\begin{tcolorbox}[title=Rappels utiles]
\begin{itemize}[leftmargin=1.2cm]
    \item Un nombre pair s'écrit sous la forme \(2n\), avec \(n\in\mathbb{Z}\).
    \item Un nombre impair s'écrit sous la forme \(2n+1\), avec \(n\in\mathbb{Z}\).
    \item Si un produit est nul, alors au moins un des facteurs est nul.
\end{itemize}
\end{tcolorbox}

\vspace{0.6em}
\textbf{Barème indicatif :} 1 pt ; 1 pt ; 1 pt.

% ========================= EXERCICE 5 =========================
\section*{Exercice 5 -- Démontrer le théorème de Pythagore \hfill /4}

\textbf{Consigne :} \emph{Démontrer le théorème de Pythagore sans schéma.}

\medskip

\textbf{Aide 1 :} vecteurs, norme d'un vecteur, produit scalaire\\
\textbf{Aide 2 :} relation de Chasles

\medskip

On considère un triangle \(ABC\) rectangle en \(A\).

On souhaite démontrer que :
\[
BC^2=AB^2+AC^2.
\]

\begin{tcolorbox}[title=Indications]
On pourra utiliser la relation de Chasles :
\[
\overrightarrow{BC}=\overrightarrow{BA}+\overrightarrow{AC}.
\]
Puis calculer la norme au carré de \(\overrightarrow{BC}\), c'est-à-dire \(\|\overrightarrow{BC}\|^2\).

On rappelle que si deux vecteurs sont perpendiculaires, alors leur produit scalaire est nul.
\end{tcolorbox}



\begin{tcolorbox}[title=Attendus]
On attend une démonstration rédigée, avec une justification claire de chaque étape : utilisation de Chasles, développement, annulation du produit scalaire, puis traduction en longueurs.
\end{tcolorbox}

\vspace{0.6em}


\vfill

\begin{center}
    \rule{0.8\textwidth}{0.6pt}\\[0.4em]
    \textit{Fin du sujet -- Toute réponse doit être justifiée avec soin.}
\end{center}

\end{document}